\section{Forced means and finite tails}\label{sec:tails}
This section supplies the bounds that hold for every positive
semidefinite Gram matrix of the corner, whatever its corrections: the
forced low means, the generating kernels of the fixed baselines, a
global bound for the pure blocks, and explicit finite tails.


\begin{lemma}[Forced complete means]\label{lem:means} Let
\(s_g=\sum_{p+q=g}\mathbf e_{pq}\) in the mixed coordinates. Then \(B\succeq0\),
and every PSD Gram \eqref{eq:corner-sectors} satisfies
\begin{equation}
 Ms_g=Bs_g=0\qquad(0\le g<m).                          \label{eq:forced-low-means}
\end{equation}

\end{lemma}
\begin{proof} The matrix $B$ is block diagonal by totals $g=p+q$. On a
complete total $g$, that is, when all $p,q\ge0$ with $p+q=g$ are
retained, put \(d=\min(p+1,g-p+1)\). Then
\[
 \sum_{r=0}^g\min(d,r+1,g-r+1)
 =\sum_{h=1}^{d}(g+3-2h)
 =d(g+2-d)=(p+1)(g-p+1),
\]
so every row of the block sums to zero. Its off-diagonal entries are
nonpositive, so it is a weighted graph Laplacian, positive semidefinite
with $s_g$ in its kernel. A block with $g\ge m$ is a principal submatrix
of the complete block given by the same formula. Hence $B\succeq0$, and
$Bs_g=0$ for $g<m$, since these totals are complete.

Now let $Q\succeq0$ represent $\calB_m$, and substitute
\(w_p^-=(t^p,0)\), \(w_p^+=(0,t^p)\). The pure wedges vanish and
$z^{\rm mix}_{pq}=t^{p+q}$, so the substituted polynomial is
$\|\sum_gt^gQ^{1/2}s_g\|^2$, with $s_g$ placed in the mixed coordinates.
Computed from the baseline, it equals $\sum_{g,h}t^{g+h}s_g^TBs_h$,
which has no term of degree below $2m$. If $Q^{1/2}s_g\ne0$ for some
$g<m$, the least such $g$ makes the coefficient of $t^{2g}$ equal to
$\|Q^{1/2}s_g\|^2>0$, a contradiction. Hence $Qs_g=0$; in particular the
mixed block satisfies $Ms_g=0$. This proves \eqref{eq:forced-low-means}. \end{proof}

Only these \emph{complete low means} are forced; the means of the
incomplete totals \(m\le g\le2m-2\) need not vanish. Symmetry and \eqref{eq:forced-low-means} imply
that \((M-B)(\xi,\xi,\zeta,\zeta)\) contains only terms \(\xi^g\zeta^h\) with both \(g,h\ge m\).
They will be bounded as finite tails.

To sum the fixed baselines, put
\[
 A=(1-\xi\zeta)(1-\eta\omega),\quad \Delta{}=(1-\xi\omega)(1-\eta\zeta),\quad
 C=(1-\xi\eta)(1-\zeta\omega),\quad \Pi=\xi\eta\zeta\omega.
\]
Absolute geometric series in the open polydisk give
\begin{equation}
\begin{aligned}
 D_\infty&=2(A^{-2}-\Delta{}^{-2}),\\
 B_\infty&=2(A^{-2}-(A\Delta{})^{-1}),\\
 Q_{0,\infty}&=D_\infty-2(\Delta{}-A)/(A\Delta{}C),\\
 L_\infty&=2\{2/A^2+1/C^2-1/(AC)-2/\Delta{}^2\}.
\end{aligned}                                                    \label{eq:generating-kernels}
\end{equation}
Appendix~\ref{app:kernels} derives these sums.

\begin{lemma}[Global bound and the repeated-node identity]\label{lem:global} For every feasible
finite averaged positive semidefinite Gram and \(|\xi|,|\eta|<1\),
\begin{equation}
 0\le P_\pm[\xi,\eta]\le
                  \frac8{(1-|\xi|^2)^2(1-|\eta|^2)^2}.      \label{eq:pure-global-bound}
\end{equation}
Writing \(\Psi=M-B\), one also has
\begin{equation}
 P_+[\xi,\bar \xi]=Q_{0,m}[\xi,\bar \xi]
                         +\Psi(\xi,\xi,\bar \xi,\bar \xi).        \label{eq:repeated-node}
\end{equation}
For \(\xi=\sqrt\rho e^{i\vartheta}\), the infinite baseline on this diagonal is
\begin{equation}
 Q_{0,\infty}[\xi,\bar \xi]=
 \frac{8\rho\sin^2\vartheta}
 {(1-\rho)^2((1-\rho)^2+4\rho\sin^2\vartheta)^2}.           \label{eq:repeated-baseline}
\end{equation}

\end{lemma}
\begin{proof} Alternation of \(E\), before using positivity, gives
\begin{equation}
 P_+[\xi,\eta]+P_-[\xi,\eta]+P_+[\xi,\bar \eta]+P_-[\xi,\bar \eta]
                        =2D_m[\xi,\eta]+2D_m[\xi,\bar \eta].     \label{eq:partial-conjugation}
\end{equation}
Every term on the left is nonnegative. The diagonal \(D_m\) is bounded
by its positive infinite sum, at most
\(2/((1-|\xi|^2)^2(1-|\eta|^2)^2)\), proving \eqref{eq:pure-global-bound}.
On the repeated slots \((\xi,\bar \xi,\bar \xi,\xi)\), \(E\) vanishes and
\[
 \Psi(\xi,\xi,\bar \xi,\bar \xi)
     =-\Theta(\xi,\bar \xi,\xi,\bar \xi)=\Theta(\xi,\bar \xi,\bar \xi,\xi),
\]
which proves the plus sign in \eqref{eq:repeated-node}. Both \eqref{eq:repeated-node} and \eqref{eq:partial-conjugation} are algebraic
identities for every averaged Gram; forced means enter only when bounding
the correction in \eqref{eq:repeated-node}. Substitution in \eqref{eq:generating-kernels}, with
\(A=C=(1-\rho)^2\) and \(\Delta{}=A+4\rho\sin^2\vartheta\), gives \eqref{eq:repeated-baseline}. \end{proof}

\begin{lemma}[Uniform finite tails]\label{lem:tails} For \(0<\epsilon\le1\), set
\begin{equation}
 \mathcal E(m,\epsilon)=2^{32}\epsilon^{-8}e^{-\epsilon m/4}.
                                                               \label{eq:tail-error}
\end{equation}
If all four slot moduli are at most \(e^{-\epsilon/2}\), then
\(|L_m-L_\infty|\le\mathcal E\). For every averaged positive semidefinite Gram, also
\(|P_+[\xi,\bar \xi]-Q_{0,\infty}[\xi,\bar \xi]|\le\mathcal E\)
when \(|\xi|\le e^{-\epsilon/2}\).

\end{lemma}
\begin{proof} Antisymmetry gives \(\operatorname{diag}(\mathcal R\Theta)=0\).
Thus the PSD matrices \(M,B\) have the same diagonal, bounded by
\(2(p+1)(q+1)\). Cauchy–Schwarz yields, with
\(W=(p+1)(q+1)(r+1)(s+1)\),
\begin{equation}
 |\Psi_{pq,rs}|\le4\sqrt W\le4W.             \label{eq:mixed-coefficients}
\end{equation}
Every expanded coefficient of \(Q_{0,\infty}\) and \(L_\infty\) has
absolute value at most \(16W\): the diagonal and minimum terms in \eqref{eq:corner-baselines}
give bounds \(2W\) for \(D\), \(2W\) for \(K\), and \(4W\) for \(B\);
even the triangle bound \(10W\) suffices for \(L_\infty\). In a discarded term at least one index
is at least \(m\). Split its radial decay in half, choose this index in
four ways, and take the product of four copies of
\(\sum_{j\ge0}(j+1)e^{-\epsilon j/4}=(1-e^{-\epsilon/4})^{-2}\).
Since \(1-e^{-\epsilon/4}\ge\epsilon/8\), the truncation error is at most
\begin{equation}
 64e^{-\epsilon m/4}(1-e^{-\epsilon/4})^{-8}
                  \le2^{30}\epsilon^{-8}e^{-\epsilon m/4}.       \label{eq:baseline-tail}
\end{equation}
The repeated correction in \eqref{eq:repeated-node} has both total degrees at least \(m\)
by Lemma~\ref{lem:means}. Using \eqref{eq:mixed-coefficients} in the same sum bounds it by
\(2^{26}\epsilon^{-8}e^{-\epsilon m/2}\). Combining these bounds proves
\eqref{eq:tail-error}, uniformly over every allowed \(E,\Theta\). \end{proof}

For \(k\ge16\), \(\epsilon=2^{-k}\), \(m=\epsilon^{-2}\), the elementary
inequality \(2^{k-2}\ge8k+40\) follows by checking \(k=16\) and doubling.
Using \(e>2\) gives
\begin{equation}
 \mathcal E\le2^{32+8k-2^{k-2}},\qquad
 \mathcal E\le\epsilon^{-2},\qquad
                    \epsilon^4\mathcal E\le\epsilon.   \label{eq:dyadic-tail-budget}
\end{equation}
The last two claims require respectively \(2^{k-2}\ge6k+32\) and
\(2^{k-2}\ge5k+32\), both consequences of the stated inequality.
